\documentclass[11pt]{article}
\usepackage[margin=0.92in]{geometry}
\usepackage[T1]{fontenc}
\usepackage[utf8]{inputenc}
\usepackage{lmodern}
\usepackage{microtype}
\usepackage{amsmath,amssymb,amsthm,mathtools,bm}
\usepackage{booktabs,array,tabularx,longtable}
\usepackage{enumitem}
\usepackage[dvipsnames]{xcolor}
\usepackage{hyperref}
\usepackage{fancyhdr}
\usepackage{titlesec}

\definecolor{TitleBlue}{HTML}{173753}
\definecolor{AccentBlue}{HTML}{2F5D7E}
\definecolor{SoftBlue}{HTML}{EEF4FA}
\definecolor{SoftGold}{HTML}{FAF6EE}
\definecolor{SoftGreen}{HTML}{EEF8F0}
\definecolor{SoftRed}{HTML}{FCEEEE}
\definecolor{RuleGray}{HTML}{D7E1EA}
\definecolor{MutedText}{HTML}{4E5A66}

\hypersetup{colorlinks=true,linkcolor=AccentBlue,urlcolor=AccentBlue,
  citecolor=AccentBlue,
  pdftitle={Type-1 N=3 LZ -- Curated bibliography for the closed-form derivation},
  pdfauthor={Type-1 N=3 Landau-Zener project}}
\pagestyle{fancy}\fancyhf{}
\fancyhead[L]{\textcolor{MutedText}{\small Curated bibliography}}
\fancyhead[R]{\textcolor{MutedText}{\small Type-1, $N=3$ Landau--Zener}}
\fancyfoot[C]{\textcolor{MutedText}{\thepage}}
\setlength{\headheight}{14pt}\renewcommand{\headrulewidth}{0.4pt}
\titleformat{\section}{\large\bfseries\color{TitleBlue}}{\thesection}{0.6em}{}
\titleformat{\subsection}{\normalsize\bfseries\color{AccentBlue}}{\thesubsection}{0.5em}{}
\setlength{\parindent}{0pt}\setlength{\parskip}{0.5em}
\setlist[itemize]{leftmargin=1.4em,itemsep=0.18em,topsep=0.18em}

\newcommand{\C}{\mathbb C}\newcommand{\R}{\mathbb R}
\newcommand{\ii}{\mathrm i}\newcommand{\dd}{\mathrm d}

\newcommand{\infobox}[2]{\begin{center}\fcolorbox{RuleGray}{#1}{\begin{minipage}{0.94\linewidth}\small #2\end{minipage}}\end{center}}

\newcommand{\bibentry}[5]{%
  \begin{center}\fcolorbox{RuleGray}{#5}{\begin{minipage}{0.94\linewidth}\small
    \textbf{#1}\par\vspace{0.3em}
    \textit{Relevant sections:} #2\par\vspace{0.2em}
    \textit{Annotation:} #3\par\vspace{0.2em}
    \textit{Priority:} #4
  \end{minipage}}\end{center}\vspace{0.4em}
}

\begin{document}

\begin{center}
{\Large\color{AccentBlue}\bfseries Curated bibliography\par}\vspace{0.45cm}
{\fontsize{19}{23}\selectfont\bfseries\color{TitleBlue}Type-1, $N=3$ Landau--Zener\par}
{\fontsize{19}{23}\selectfont\bfseries\color{TitleBlue}closed-form derivation\par}
\vspace{0.30cm}
{\large\color{MutedText}Reading roadmap for the deep literature review\par}
\end{center}

\infobox{SoftBlue}{%
\textbf{Purpose.} This document is a curated, annotated bibliography for
the open closed-form derivation of the Type-1, $N=3$ ungauged multistate
Landau--Zener transition matrix. Phase I delivered two Brundobler--Elser
extreme-survival entries exactly. Phase II tested three exact-WKB
literature routes and all failed empirically. Phase III narrowed the open
problem to: identify a $\mathrm U(3)$-valued function
$D(\epsilon, \gamma, a)$ that produces the remaining two transcendental
entries of $P$.

Each reference below is annotated with what it gives us, what to
extract on a careful read, and a priority tier:
\textbf{P1} (must read first; resolves a Phase II/III gap),
\textbf{P2} (secondary; framework context),
\textbf{P3} (supplementary; useful for write-up). A 5-pass reading
roadmap appears at the end.}

\tableofcontents

\vspace{0.5em}

% ===========================================================================
\section{Spectral networks (Gaiotto--Moore--Neitzke)}

The foundational framework for the WKB connection matrix as a geometric
object on a Riemann surface. The Type-1, $N=3$ spectral curve
$\Sigma_Y : \mu^2 = Q_4 L_H^2/p^2$ (genus 1, $b_1 = 2$) sits squarely
inside this framework.

\bibentry{Gaiotto, Moore, Neitzke. \textit{Spectral networks}.
arXiv:1204.4824 (2012). Annales Henri Poincar\'e 14 (2013) 1643--1731.}
{\S\S5--6 (wall structure, BPS counting), \S10 (non-abelianisation map).}
{The geometric language of spectral networks: BPS trajectories
emanating from turning points, walls of marginal stability, and the
non-abelianisation map from non-abelian flat connections to flat
$\mathrm U(1)$ connections on the spectral cover. \textbf{What it
gives us:} the topological skeleton on which the inter-window
connection matrix lives; the precise definition of the $\omega$-pair
$K_\pm(\gamma) = I + \mu_\gamma E_{ij}, \, I + (1/\bar\mu_\gamma) E_{ji}$.
\textbf{What to extract:} \S5.3 the sliding-flag construction;
\S6.2 the $K$-wall factor and its reciprocal-$\mu$ pair; \S10.2 the
abelianisation map; \S10.3 the ordered chamber product.}
{P1}{SoftBlue}

\bibentry{Gaiotto, Moore, Neitzke. \textit{Spectral networks and snakes}.
arXiv:1209.0866 (2012). Annales Henri Poincar\'e 15 (2014) 61--141.}
{\S3 (abelianisation map on snake diagrams).}
{Cluster-coordinate refinement of GMN 2012. Snake diagrams are the
combinatorial encoding of the spectral network on a punctured surface
useful for computing the connection matrix entry-by-entry.
\textbf{What it gives us:} a concrete worked example of non-abelianisation
in low rank with explicit formulas. \textbf{What to extract:} \S3.2
the cluster-coordinate decomposition of the connection matrix; relate
to Iwaki--Nakanishi (\S3 below).}
{P2}{SoftGold}

% ===========================================================================
\section{Non-abelianisation and Hitchin systems (Hollands--Neitzke)}

The direct reference family for the project's setting. Hollands--Neitzke
extend GMN to explicit connection-matrix formulas for rank-2 and rank-3
Hitchin systems, including the genus-1 case that matches our
$\Sigma_Y$ exactly.

\bibentry{Hollands, Neitzke. \textit{Spectral networks and Fenchel--Nielsen
coordinates}. arXiv:1312.2979 / 1402.1131 (2016). Lett. Math. Phys.
106 (2016) 811--877.}
{\S\S4 (explicit rank-2 abelianisation), \S4.4 (the $\omega$-pair to
$N^{\mathrm{proper}}$ collapse).}
{The proof we need most: that for a rank-2 spectral network with the
canonical BPS-ray choice, the inner ω-pair product
$K_-(\gamma) \cdot \Phi_{\mathrm{mid}} \cdot K_+(\gamma)$ collapses
\emph{exactly} to the Joye--Kunz--Pfister Stokes-phased
$N^{\mathrm{proper}}$ block. \textbf{What it gives us:} the algebraic
identity that ratifies the expected Route A closed form
$U_{HN} = \Phi_R N_B^{\mathrm{proper}} \Phi_{AB} N_A^{\mathrm{proper}} \Phi_L$.
\textbf{What to extract:} \S4.4 Lemma giving the collapse;
\S4.2 the explicit ${\rm SL}(2,\mathbb C)$-valued formula.}
{P1}{SoftBlue}

\bibentry{Hollands, Neitzke. \textit{BPS states in spectral networks of
arbitrary rank}. arXiv:1611.04317 (2016).}
{\S2 (rank-$N$ generalisation), \S5 (genus-$\ge 1$ examples).}
{Rank-3 / higher-rank generalisation that includes our $N=3$ case
without reduction to rank 2. \textbf{What it gives us:} confirmation
that the rank-3 ω-pair construction works on genus-1 Hitchin systems
and that the BPS-state count is fixed by $b_1$ of the cover.
\textbf{What to extract:} \S2.3 the rank-$N$ wall-jump formula; the
worked $\mathrm{SL}(3)$ example for spectral curves with $b_1 = 2$.}
{P2}{SoftGold}

\bibentry{Hollands, Neitzke. \textit{Exact WKB and abelianization for the
$T_3$ equation}. arXiv:1906.04271 (2019). Commun. Math. Phys.
388 (2021) 1571.}
{\S\S5--6 (chamber-structure proof for genus-1 covers with $b_1 = 2$).}
{\textbf{The direct genus-1 reference for our setting.} The $T_3$
equation is a third-order linear ODE whose spectral cover has exactly
the topology of our $\Sigma_Y$ (genus 1, $b_1 = 2$). \textbf{What it
gives us:} the proof that the GMN ordered product at the canonical
phase collapses to two $N^{\mathrm{proper}}$ factors plus a diagonal
inter-window factor $\Phi_{AB}$, with no extra inter-window BPS state
in the relevant chamber. This is the load-bearing theorem for
Route A's expected closed form.
\textbf{What to extract:} \S5 chamber-structure analysis; \S6 the
final connection-matrix formula in the relevant chamber; convention
table for $\theta$-phase and ω-pair side choice.}
{P1}{SoftBlue}

% ===========================================================================
\section{Iwaki--Nakanishi cluster algebras}

The cluster-algebra dual of GMN: Voros symbols become cluster
$Y$-variables and Stokes phenomena become cluster mutations. Provides
an alternative vocabulary for the same connection matrix.

\bibentry{Iwaki, Nakanishi. \textit{Exact WKB analysis and cluster
algebras}. arXiv:1401.7094 (2014). J. Phys. A 47 (2014) 474009.}
{\S\S3--4 (cluster mutation as Stokes phenomenon), \S5 (the connection
matrix formula as mutation product).}
{The cluster-algebraic vocabulary for exact WKB. Each mutation
contributes a Faddeev quantum dilogarithm $\Phi_b$ factor; the
connection matrix is a finite product of these along a cluster path.
\textbf{What it gives us:} Route B's framework; explicit formula
for the mutation product that produces a $\mathrm U(N)$ connection
matrix from $N+1$ cluster variables. \textbf{What to extract:}
\S3.2 mutation rule; \S4 the $T$-system / $Y$-system identity that
our Sinitsyn--Chernyak Track 3 result corresponds to.}
{P1}{SoftBlue}

\bibentry{Iwaki, Nakanishi. \textit{Voros coefficients of the Schr\"odinger
equation}. arXiv:1409.4641 (2014). J. Phys. A 47 (2014) 474012.}
{\S\S2--3 (Voros symbols as cluster coordinates).}
{Companion paper to 1401.7094, focused on the Schr\"odinger pencil
$\hbar^2 \psi'' + Q(u) \psi = 0$. \textbf{What it gives us:} explicit
Voros-symbol formulas for the rank-2 reduction of our quotient-ring
scalar equation. \textbf{What to extract:} the formula relating
$\oint \sqrt Q\,\dd u$ to the cluster $Y$-variable, and the Faddeev
$\Phi_b$ argument at finite $b$.}
{P2}{SoftGold}

% ===========================================================================
\section{Quantum dilogarithm (Faddeev)}

The fundamental special function of the cluster-mutation route. Plays
the same role for cluster algebras that $\Gamma$ plays for the
Joye--Kunz--Pfister Stokes block.

\bibentry{Faddeev. \textit{Discrete Heisenberg-Weyl group and modular
group}. Lett. Math. Phys. 34 (1995) 249--254.}
{Eq. (5) (the $\Phi_b$ definition); \S2 (the five-term relation).}
{Original definition of the non-compact / modular quantum dilogarithm
$\Phi_b(z)$ as a contour integral or infinite product. \textbf{What
it gives us:} the function our finite-$b$ Faddeev refactor needs.
\textbf{What to extract:} the integral representation (better
behaved than the infinite product for large $|z|$); the
$b \to b^{-1}$ duality.}
{P1}{SoftBlue}

\bibentry{Faddeev, Kashaev. \textit{Quantum dilogarithm}. Mod. Phys. Lett.
A 9 (1994) 427--434.}
{\S\S2--3 (basic properties; semiclassical limit).}
{The original integral representation of $\Phi_b$. \textbf{What it
gives us:} the $b \to 0$ classical limit (recovers
$\mathrm{Li}_2$); the connection to $\Gamma$-functions in the
appropriate limit. \textbf{What to extract:} the precise small-$b$
asymptotic expansion.}
{P2}{SoftGold}

% ===========================================================================
\section{Aoki--Kawai--Takei exact WKB / DDP}

The framework closest to Phase II's Route C: discontinuity formulas
for exact-WKB connection matrices at simple and coalescing turning
points.

\bibentry{Aoki, Kawai, Takei (and collaborators). \textit{Exact WKB
analysis} (review). E.g. RIMS preprint series and Kyoto lecture
notes, ~1992--2010.}
{\S\S4--5 (DDP at simple and coalescing turning points; Whittaker /
Bessel matching).}
{The exact-WKB connection-matrix programme. \textbf{What it gives
us:} the proper formula for the connection matrix at two coalescing
turning points (modified-Bessel index 1/3 at argument $2|V|/3$).
\textbf{What to extract:} the Wronskian-normalised form of the
Bessel block (Phase III prototype used a heuristic
$\sin(\arg V)$ which was wrong); the precise sheet selection for
the inter-window Voros symbol.}
{P1}{SoftBlue}

\bibentry{Delabaere, Dillinger, Pham. \textit{R\'esurgence de Voros et
p\'eriodes des courbes hyperelliptiques}. Ann. Inst. Fourier 43
(1993) 163--199.}
{Theorems II.2, II.3 (the discontinuity formula at a simple turning
point).}
{The original DDP discontinuity formula: the analytic continuation
of a WKB symbol around a turning point picks up a controlled
$\Gamma$-function-valued jump. \textbf{What it gives us:} the
elementary connection block we already implement in
\texttt{assay.ddp.j\_ddp}; the framework that AKT extends to
coalescing turning points. \textbf{What to extract:} Theorem II.3
(the universal form of the discontinuity); the residue computation
at the turning point.}
{P1}{SoftBlue}

% ===========================================================================
\section{Multistate Landau--Zener -- historical core}

\bibentry{Brundobler, Elser. \textit{S-matrix of multistate
Landau--Zener problem}. J. Phys. A 26 (1993) 1211--1227.}
{Throughout; especially the extreme-level survival formula.}
{The original derivation of the Brundobler--Elser product:
$P[\text{extreme level survives}] = \prod q_{ij}$ with $q_{ij} =
\exp(-2\pi \Gamma_{ij})$. \textbf{What it gives us:} the two
exact-elementary entries of the Type-1 $N=3$ matrix $P$ (Phase I
deliverable). \textbf{What to extract:} the topological argument
that fixes the two extreme corners independently of all integrability
data; the implicit ``no-self-interaction'' assumption.}
{P1}{SoftBlue}

\bibentry{Demkov, Osherov. \textit{Stationary and nonstationary problems
in quantum mechanics that can be solved by the method of
factorization}. Sov. Phys. JETP 24 (1967) 916.}
{Sections on the equal-slope special case.}
{The original three-state LZ with one ``selector'' level crossing
the other two. \textbf{What it gives us:} historical context for
the Type-1 geometry (Demkov--Osherov is the symmetric special case;
Type-1 is the generic version). \textbf{What to extract:} the
sequential-LZ-product structure that DOES hold in the
Demkov--Osherov special case; the precise extension to the generic
Type-1 is what Phase III is solving.}
{P2}{SoftGold}

\bibentry{Sinitsyn. \textit{Multistate Landau--Zener problem}.
arXiv:cond-mat/0212017 (2002) and successors; recent review in
New J. Phys. 19 (2017) 075008 with Chernyak.}
{\S\S2--3 (integrability conditions, hidden symmetry).}
{The integrable hierarchy of multistate LZ, including the
Sinitsyn--Chernyak commuting-triple condition $[H, H'] = 0$.
\textbf{What it gives us:} Track 3 verified this commuting structure
to $7\times 10^{-15}$; the hierarchy pins the two extreme diagonals
but leaves the middle survival open. \textbf{What to extract:} the
identification of which $P$ entries are integrability-fixed vs.
genuinely transcendental --- and exactly how many free parameters
remain.}
{P1}{SoftBlue}

% ===========================================================================
\section{Two-level LZ with Stokes-phase corrections}

\bibentry{Joye, Kunz, Pfister. \textit{Exponential estimates in
adiabatic quantum evolution}. Ann. Inst. H. Poincar\'e 54 (1991)
343--378. Companion: J. Math. Phys. 32 (1991) 3328.}
{\S\S2--3 (the rigorous $\varphi_S(\delta)$ formula).}
{The mathematically rigorous derivation of the Stokes phase
$\varphi_S(\delta) = \delta(\log\delta - 1) + \pi/4 +
\arg\Gamma(1 - \ii\delta)$ for the two-level LZ amplitude block.
\textbf{What it gives us:} the function we implement in
\texttt{assay.lz\_stokes\_phased}; the closed-form expression we
expect to see emerge from the GMN ω-pair construction.
\textbf{What to extract:} the proof that this is the unique
non-perturbative correction to the elementary BE amplitude in the
2-level case; the convention for the sign of $\varphi_S$.}
{P1}{SoftBlue}

\bibentry{Berry. \textit{Histories of adiabatic quantum transitions}.
Proc. R. Soc. A 429 (1990) 61--72.}
{\S\S3--4 (Stokes lines and the histories interpretation).}
{Berry's "matching" interpretation of the LZ Stokes phenomenon.
\textbf{What it gives us:} physical context for $\varphi_S$ as a
matching condition across Stokes lines, useful for explaining the
result. \textbf{What to extract:} the geometric picture of the
Stokes lines in the complex time plane and how they connect to
the GMN spectral network.}
{P2}{SoftGold}

% ===========================================================================
\section{Hitchin integrable systems (foundational)}

\bibentry{Hitchin. \textit{Stable bundles and integrable systems}.
Duke Math. J. 54 (1987) 91--114.}
{\S\S4--6 (definition; rank-$N$ case; genus-1 special case).}
{The foundational paper defining Hitchin systems. \textbf{What it
gives us:} the rigorous identification of our commuting triple
$(H_x, Y, R)$ from Track 3 as a Hitchin system on a punctured
genus-0 base with rank-3 spectral cover (which IS our $\Sigma_Y$
with $b_1 = 2$). \textbf{What to extract:} \S6 the genus-1 case;
the symplectic structure that the Hollands--Neitzke abelianisation
preserves.}
{P2}{SoftGold}

\bibentry{Hitchin. \textit{The self-duality equations on a Riemann
surface}. Proc. London Math. Soc. 55 (1987) 59--126.}
{\S\S5--6 (rank-$N$ moduli space; the integrable hierarchy).}
{Companion paper to Duke 1987; constructs the full Hitchin moduli
space. \textbf{What it gives us:} the moduli-space perspective on
why our two-parameter open content sits in a 6-dimensional U(3)
chamber. \textbf{What to extract:} the parameter count for the
rank-3 / genus-1 Hitchin moduli (matches Phase III's 6-after-gauge
count of $D$).}
{P3}{SoftGreen}

% ===========================================================================
\section{Cluster coordinates (Fock--Goncharov, Allegretti)}

\bibentry{Fock, Goncharov. \textit{Cluster ensembles, quantization and
the dilogarithm}. arXiv:math/0311149 (2003). Ann. Sci. \'EN.S.
42 (2009) 865.}
{\S\S2--3 (cluster ensembles; $X$- and $A$-coordinates).}
{The general framework of cluster ensembles. \textbf{What it gives
us:} the rigorous structure that contains Iwaki--Nakanishi as a
special case; the duality between Voros symbols (X-coordinates)
and shear coordinates (A-coordinates). \textbf{What to extract:}
the mutation rules and the dilogarithm identities that constrain
the cluster product.}
{P2}{SoftGold}

\bibentry{Allegretti. \textit{Voros symbols as cluster coordinates}.
arXiv:1810.09433 (2018). Int. Math. Res. Not. (2019).}
{Throughout; especially the cluster-coordinate identification.}
{Explicit identification of WKB Voros symbols with Fock--Goncharov
cluster coordinates. \textbf{What it gives us:} the bridge that
makes Iwaki--Nakanishi rigorous; the proof that the cluster
$Y$-variable equals $\exp(\oint \sqrt Q\, \dd u)$ along the
appropriate cycle. \textbf{What to extract:} the precise dictionary
between WKB and cluster vocabulary for our genus-1 case.}
{P1}{SoftBlue}

% ===========================================================================
\section{Special functions and reference handbooks}

\bibentry{Olver, Lozier, Boisvert, Clark (eds.). \textit{NIST Digital
Library of Mathematical Functions}. \href{https://dlmf.nist.gov}{dlmf.nist.gov},
ongoing. Companion: Abramowitz \& Stegun, \textit{Handbook of
Mathematical Functions} (1964), Ch. 9--10.}
{Ch. 9--10 (modified Bessel functions); Ch. 5--6 ($\Gamma$ and related
functions).}
{Reference for the modified Bessel functions $I_\nu, K_\nu$ at
$\nu = 1/3$ needed for the AKT block, and for the $\Gamma$-function
identities used throughout. \textbf{What it gives us:} numerical
values, asymptotic expansions, Wronskian relations.
\textbf{What to extract:} the Wronskian
$I_\nu K'_\nu - I'_\nu K_\nu = -1/z$ (DLMF 10.28.2) which fixes the
unitary normalisation of the AKT block.}
{P3}{SoftGreen}

% ===========================================================================
\newpage
\section*{Reading roadmap (5 passes)}
\addcontentsline{toc}{section}{Reading roadmap (5 passes)}

\infobox{SoftGold}{The references above are organised by framework. The
order below is the recommended \emph{reading} order, driven by which
Phase II/III gap each pass closes. For each item the annotation tells
you exactly what to extract --- this is not an open-ended read.}

\subsection*{Pass 1 -- Route A gap fix (\textbf{1--2 days})}

\textbf{Goal.} Extract the ω-pair side-tracking convention that fixes
the Phase III Route A prototype's catastrophic overflow
(residual $10^{23}$ to $10^{193}$ from naive
$I + \mu E_{ij}$ insertions with $|\mu| \sim e^{50}$).

\textbf{Read.}
\begin{itemize}
 \item Hollands--Neitzke arXiv:1906.04271 \S\S5--6 (the genus-1
 chamber-structure proof for $T_3$, matching our $\Sigma_Y$ with
 $b_1 = 2$).
 \item Gaiotto--Moore--Neitzke arXiv:1204.4824 \S5 (the
 $K_\pm(\gamma)$ pair definition) and \S10.3 (the ordered chamber
 product).
\end{itemize}

\textbf{Extract.} The exact algebraic identity
$K_-(\gamma) \cdot \Phi_{\mathrm{mid}} \cdot K_+(\gamma) =
N_X^{\mathrm{proper}}$ for an inner ω-pair, where
$N_X^{\mathrm{proper}}$ is the Joye--Kunz--Pfister Stokes-phased block.
This is the load-bearing identity for Route A's expected closed form
$U_{HN} = \Phi_R \cdot N_B^{\mathrm{proper}} \cdot \Phi_{AB} \cdot
N_A^{\mathrm{proper}} \cdot \Phi_L$.

\subsection*{Pass 2 -- Route B gap fix (\textbf{2--3 days})}

\textbf{Goal.} Extract the log-space cluster-mutation product and the
finite-$b$ Faddeev kernel that would fix Phase II Route B's Voros-$Y$
overflow.

\textbf{Read.}
\begin{itemize}
 \item Iwaki--Nakanishi arXiv:1401.7094 \S\S3--4 (the mutation rule and
 the cluster product).
 \item Faddeev Lett. Math. Phys. 34 (1995) eq. (5) and \S2 (the
 $\Phi_b$ definition and five-term relation).
\end{itemize}

\textbf{Extract.} The explicit mutation formula
$\mu_k(Y_j) = Y_j (1 + Y_k^{\mp 1})^{B_{jk}}$ with the $\Phi_b$
contribution per mutation; the convention for finite $b$ at our
$\epsilon = 1$ semiclassical limit. (Route B was not implemented in
Phase III but is architecturally sufficient if the Phase II
combinatorics is corrected.)

\subsection*{Pass 3 -- Route C gap fix (\textbf{2--3 days})}

\textbf{Goal.} Extract the correct inter-window Voros symbol (sheet
selection) and the Bessel-index-$1/3$ connection block. The Phase III
Route C prototype used a heuristic $\sin(\arg V)$ for the mixing
magnitude and integrated on the wrong sheet
($V_{\mathrm{new}} = 2 V_{\mathrm{old}}$ across all samples).

\textbf{Read.}
\begin{itemize}
 \item AKT review (Kyoto / RIMS lecture notes) \S\S4--5.
 \item Delabaere--Dillinger--Pham Ann. Inst. Fourier 43 (1993)
 \S III (the discontinuity formula).
\end{itemize}

\textbf{Extract.} The Wronskian-normalised Bessel block at index
$\nu = 1/3$ and argument $2 |V| / 3$; the precise sheet on which
$V_{AB}$ must be integrated (almost certainly the spectator sheet on
which the two windows ``join,'' not the exterior sheet $\lambda_0$
the Phase III prototype used).

\subsection*{Pass 4 -- Context for write-up (\textbf{1 day})}

\textbf{Goal.} Frame the eventual result in the established LZ
literature.

\textbf{Read.}
\begin{itemize}
 \item Brundobler--Elser J. Phys. A 26 (1993) 1211.
 \item Joye--Kunz--Pfister Ann. IHP 54 (1991).
 \item Berry Proc. R. Soc. A 429 (1990) 61.
\end{itemize}

\textbf{Extract.} The standard convention for the LZ Stokes phase
sign; the matching condition across the Stokes lines (for the
physical interpretation in the final brief).

\subsection*{Pass 5 -- Depth (\textbf{ongoing})}

\textbf{Goal.} Ensure rigorous embedding of the result in the
cluster/Hitchin framework for the eventual paper.

\textbf{Read.}
\begin{itemize}
 \item Allegretti arXiv:1810.09433 (Voros--cluster dictionary).
 \item Fock--Goncharov arXiv:math/0311149 (cluster ensembles).
 \item Hitchin Duke Math. J. 54 (1987) 91 (foundational).
\end{itemize}

\textbf{Extract.} The mapping between our Type-1 $N=3$ system and
the standard Hitchin moduli; the parameter count
(6-after-gauge $\mathrm U(3)$) as a Hitchin-moduli statement.

\vspace{0.5em}

\infobox{SoftBlue}{\textbf{Total Phase III literature budget:} ~10 days
of focused reading. Pass 1 alone (~2 days) should suffice to unblock
Route A's ω-pair implementation. Passes 2--3 are insurance against
Route A failure. Passes 4--5 are write-up depth.}

\section*{Cross-reference index}

For navigation, here is the mapping between Phase II/III gaps and the
references that close them:

\begin{center}\small
\begin{tabular}{p{0.30\linewidth} p{0.30\linewidth} p{0.32\linewidth}}
\toprule
Gap & Primary reference & Secondary \\
\midrule
Route A ω-pair side-tracking & Hollands--Neitzke 1906.04271 \S\S5--6 & GMN 1204.4824 \S5 \\
Route A genus-1 collapse to $N^{\mathrm{proper}}$ & Hollands--Neitzke 1402.1131 \S4.4 & --- \\
Route B log-space mutation & Iwaki--Nakanishi 1401.7094 \S3 & Faddeev LMP 34 (1995) \\
Route B finite-$b$ Faddeev kernel & Faddeev LMP 34 (1995) & Faddeev--Kashaev MPL A 9 \\
Route C Bessel block normalisation & AKT review \S5 & DLMF 10.28 \\
Route C inter-window $V_{AB}$ sheet & AKT review \S4 & DDP 1993 \\
Stokes phase $\varphi_S(\delta)$ & Joye--Kunz--Pfister 1991 & Berry 1990 \\
BE extreme-survival formula & Brundobler--Elser 1993 & --- \\
Integrability / commuting triple & Sinitsyn--Chernyak NJP 19 (2017) & Sinitsyn 2002 \\
Cluster coordinates & Allegretti 2018 & Fock--Goncharov 2003 \\
Hitchin moduli context & Hitchin Duke 1987 & Hitchin LMS 1987 \\
\bottomrule
\end{tabular}
\end{center}

\section*{Project context references (already shipped)}

For self-consistency, the Phase II/III status is documented in:

\begin{itemize}
 \item \texttt{window\_period\_BE\_open\_problem\_brief.pdf} ---
 the original open-problem statement.
 \item \texttt{six\_track\_synthesis\_brief.pdf} --- Phase I synthesis
 (Tracks 1--6).
 \item \texttt{phase2\_findings\_brief.pdf} (11 pages) --- the
 three-route Phase II programme and validation failure analysis.
 \item \texttt{phase3\_status\_addendum.pdf} (4 pages) --- the
 Phase III Route D negative finding (U(3) richness intrinsic).
 \item \texttt{notes/track\{1..6\}\_*.pdf} --- per-track Phase I notes.
\end{itemize}

\vspace{0.5em}

\noindent\hfill\textcolor{MutedText}{\emph{End of curated bibliography.
Total references: 18 across 9 frameworks. Primary reading: 4 P1
references covering Routes A and C plus the Stokes-phase foundation.}}\hfill\null

\end{document}
